\section{Síntesis y ampliación}

Esta sesión construye una visión mínima y completa del Transformer:

\begin{enumerate}
  \item \textbf{La historia no es una ruptura}: RNN/seq2seq dejan expuesto el fixed-vector bottleneck, la attention temprana ofrece lectura bajo demanda y la self-attention la convierte en la capa principal de secuencia;
  \item \textbf{La attention es routing}: la query expresa la necesidad, la key ofrece el índice de coincidencia y el value aporta el contenido que se transmite;
  \item \textbf{El espacio de diseño ya existía}: soft/hard y global/local controlan, respectivamente, la selección diferenciable y el alcance del acceso;
  \item \textbf{La self-attention acorta las rutas}: los token pueden interactuar directamente y entrenarse en paralelo, pero la full score matrix tiene un costo cuadrático;
  \item \textbf{Un block completo no es solo attention}: position, mask, residual, normalization y FFN resuelven, respectivamente, el orden, la visibilidad y la optimización en profundidad;
  \item \textbf{Las familias se definen por las fronteras de información}: encoder-only produce representaciones bidireccionales, decoder-only genera de forma causal y encoder-decoder procesa seq2seq condicional;
  \item \textbf{GPT-3 cambió la interfaz de las tareas}: el preentrenamiento decoder-only permite expresar las tareas mediante prompt y example;
  \item \textbf{BERT reforzó la interfaz de representación}: el MLM encoder-only aprende contextual representation bidireccional;
  \item \textbf{Pasar a otros dominios no es copiar}: cada dominio debe redefinir token, position, mask, objective y evaluation;
  \item \textbf{La arquitectura no resuelve por sí sola los problemas futuros}: memory, efficiency, alignment y deployment reliability siguen requiriendo un diseño propio.
\end{enumerate}

\begin{importantbox}{Ejercicio de cálculo a mano y de diseño}
Dados $Q,K,V$ para tres token, calcula a mano el score, el scaling, la softmax y el output; después agrega por separado una padding mask y una causal mask, y explica qué pesos deben valer cero. Luego elige uno de estos casos: imágenes, RL trajectory o secuencias de proteínas; define token, position, attention pattern, objective y metric, y explica por qué usarías encoder-only, decoder-only o encoder-decoder.
\end{importantbox}

\begin{warningbox}{Autoevaluación final}
Si solo puedes repetir «el Transformer usa attention», todavía no dominas la interfaz. Debes poder explicar en qué difieren los orígenes de la basic attention y de la self-attention, cómo la mask define las fronteras de información, por qué GPT y BERT comparten el block pero sirven a tareas distintas, y si $O(n^2)$ es una propiedad de la estructura matemática o un cuello de botella de una implementación concreta.
\end{warningbox}

\subsection{Lecturas adicionales}

{\small
\begin{itemize}[nosep,leftmargin=*]
  \item \href{https://arxiv.org/abs/1409.0473}{\textbf{Bahdanau et al. (2014)}}: additive attention en encoder-decoder y alignment.
  \item \href{https://arxiv.org/abs/1508.04025}{\textbf{Luong et al. (2015)}}: global/local attention y multiplicative scoring.
  \item \href{https://arxiv.org/abs/1706.03762}{\textbf{Vaswani et al. (2017)}}: scaled dot-product, multi-head y el Transformer encoder-decoder original.
  \item \href{https://arxiv.org/abs/2005.14165}{\textbf{Brown et al. (2020)}}: GPT-3 y la evaluación few-shot/in-context.
  \item \href{https://arxiv.org/abs/1810.04805}{\textbf{Devlin et al. (2018)}}: BERT, MLM, NSP y la interfaz de fine-tuning.
  \item \href{https://web.stanford.edu/class/cs25/past/cs25-v1/}{\textbf{Stanford CS25 V1}}: la speaker series original y el programa del curso.
\end{itemize}
}

\end{document}
